\documentclass[11pt,a4paper]{article}

% Encoding and language
\usepackage[utf8]{inputenc}
\usepackage[T1]{fontenc}
\usepackage[spanish]{babel}

% Fonts
\usepackage{lmodern}
\usepackage{microtype}

% Page layout
\usepackage[margin=2.5cm]{geometry}
\usepackage{setspace}
\onehalfspacing

% Colors and boxes
\usepackage[dvipsnames]{xcolor}
\usepackage{tcolorbox}
\tcbuselibrary{skins,breakable}

% Lists
\usepackage{enumitem}

% Hyperlinks
\usepackage[colorlinks=true,linkcolor=MidnightBlue,urlcolor=MidnightBlue]{hyperref}

% Custom box style for form fields
\newtcolorbox{campo}[1][]{
    colback=gray!5,
    colframe=gray!40,
    fonttitle=\bfseries,
    title=#1,
    breakable,
    boxrule=0.5pt,
    arc=2pt,
    left=8pt,
    right=8pt,
    top=6pt,
    bottom=6pt
}

\begin{document}

\begin{center}
{\LARGE \textbf{Propuesta de Trabajo Final de Máster}}\\[0.3cm]
{\large Campos para el formulario de registro}\\[0.5cm]
\hrule
\end{center}

\vspace{0.8cm}

\begin{campo}[Título provisional]
Traducción automática Inga-Español mediante adaptación de modelos de lenguaje con RAG para la preservación digital de una lengua indígena colombiana
\end{campo}

\vspace{0.5cm}

\begin{campo}[Descripción y justificación del trabajo a desarrollar]
La lengua Inga, de la familia quechua, es hablada por 18.000 personas en el Putumayo, Colombia. Actualmente no existe ningún traductor automático que la contemple: Google Translate y DeepL la excluyen, y los recursos digitales disponibles son escasos y sin integración computacional. Esta ausencia tecnológica profundiza la brecha digital de las lenguas indígenas.

Este trabajo propone desarrollar un sistema de traducción Inga-Español combinando dos enfoques: (1) adaptación del modelo NLLB-200 de Meta mediante fine-tuning con LoRA, aprovechando que ya incluye Quechua Ayacucho como base de transferencia; y (2) uso de Claude Opus 4.5 con RAG (Retrieval-Augmented Generation) para inyectar conocimiento lingüístico del Inga en el proceso de traducción.

El proyecto nace desde dentro de la comunidad: uno de los integrantes es indígena Inga, hablante nativo y residente en Mocoa, lo que garantiza acceso directo a hablantes para construcción y validación del corpus, además de comprensión del contexto cultural.
\end{campo}

\vspace{0.5cm}

\begin{campo}[Objetivos e impacto que se espera conseguir con el desarrollo del trabajo]
\textbf{Objetivo general:} Desarrollar y evaluar un sistema de traducción Inga-Español que combine modelos de código abierto con fine-tuning y modelos de frontera con RAG, contribuyendo a la preservación digital de la lengua Inga.

\textbf{Objetivos específicos:}
\begin{enumerate}[nosep,leftmargin=*]
    \item Construir un corpus paralelo Inga-Español (meta: 5.000-10.000 pares) con participación comunitaria.
    \item Implementar una base de conocimiento lingüístico (diccionario, gramática, ejemplos) indexada para RAG.
    \item Adaptar NLLB-200 mediante LoRA y comparar con MADLAD-400 y TranslateGemma.
    \item Implementar pipeline de traducción con Claude Opus 4.5 + RAG.
    \item Evaluar con métricas automáticas (BLEU, chrF++, BERTScore) y validación con hablantes nativos.
\end{enumerate}

\textbf{Impacto esperado:}
\begin{itemize}[nosep,leftmargin=*]
    \item \textit{Social:} Herramienta para maestros bilingües, gestores culturales y organizaciones indígenas del Putumayo.
    \item \textit{Académico:} Evidencia comparativa sobre fine-tuning vs. RAG para lenguas de muy bajos recursos.
    \item \textit{Comunitario:} Corpus y recursos digitales liberados como bienes comunes para la comunidad Inga.
\end{itemize}
\end{campo}

\vspace{0.5cm}

\begin{campo}[Metodología, tecnologías o técnicas previstas para abordar el desarrollo del trabajo]
\textbf{Metodología:} Investigación aplicada con enfoque experimental comparativo en cinco fases: (1) construcción del corpus paralelo, (2) base de conocimiento para RAG, (3) fine-tuning de modelos locales, (4) pipeline Claude + RAG, (5) evaluación comparativa.

\textbf{Modelos:}
\begin{itemize}[nosep,leftmargin=*]
    \item Local principal: NLLB-200-3.3B (Meta) --- ya contiene Quechua, dominante en AmericasNLP 2021-2025.
    \item Alternativos: MADLAD-400-10B, TranslateGemma-12B.
    \item API: Claude Opus 4.5 (Anthropic) --- demostrada eficiencia en lenguas de bajos recursos.
\end{itemize}

\textbf{Técnicas:}
\begin{itemize}[nosep,leftmargin=*]
    \item Fine-tuning con LoRA/QLoRA (PEFT de HuggingFace).
    \item RAG con índices separados (léxico, gramatical, ejemplos) usando FAISS/ChromaDB y LlamaIndex.
    \item Augmentación de datos con back-translation y generación sintética.
\end{itemize}

\textbf{Infraestructura:}
\begin{itemize}[nosep,leftmargin=*]
    \item Hardware local: Apple M4 Max, 128 GB RAM (modelos hasta 40B en precisión completa).
    \item API: Anthropic Claude Opus 4.5.
\end{itemize}

\textbf{Equipo:} Tres integrantes distribuidos geográficamente:
\begin{itemize}[nosep,leftmargin=*]
    \item Mocoa (Putumayo): Corpus y validación comunitaria.
    \item Cali: RAG y orquestación del sistema.
    \item Cúcuta: Fine-tuning y evaluación.
\end{itemize}
\end{campo}

\end{document}
